\section{树递归神经网络}

\subsection{语言的递归结构}

本课的后半部分转向一种完全不同的建模思路——\textbf{树递归神经网络（Tree Recursive Neural Networks）}。这是 Manning 和学生在 Stanford 从 2010 年开始发展的一系列工作。

\begin{figure}[H]
\centering
\includegraphics[width=0.95\textwidth]{slides-images/slide-035.jpg}
\caption{语言的递归结构：名词短语可以嵌套包含名词短语，形成递归的层次结构\protect\footnotemark}
\end{figure}
\footnotetext{Slides 第36页。}

其动机来自语言学中对人类语言结构的基本观察：

\begin{importantbox}{语言的递归性}
人类语言的一个核心特性是\textbf{递归结构}——同样的句法结构可以嵌套在自身内部。例如：

``\textit{[The person standing next to [the man from [the company that purchased [the firm that you used to work at]]]]}.''

这个句子包含了四层嵌套的名词短语，每一层都有相同的句法结构。这种无限嵌套的能力是人类语言区别于其他动物交流系统的关键特征。
\end{importantbox}

\begin{figure}[H]
\centering
\includegraphics[width=0.95\textwidth]{slides-images/slide-037.jpg}
\caption{Penn Treebank 中的句法分析树示例：展示了动词短语的多层嵌套结构\protect\footnotemark}
\end{figure}
\footnotetext{Slides 第38页。}

\subsection{组合性原则}

树递归神经网络的理论基础是语言学和语言哲学中的\textbf{组合性原则（Principle of Compositionality）}：

\begin{knowledgebox}{Frege 的组合性原则}
一个短语或句子的意义由其组成词的意义和组合它们的规则共同决定。

如果我们有词向量来表示词的意义，那么我们应该能够根据句法结构，通过某种组合函数，递归地构建短语和句子的意义表示。例如，``the country of my birth'' 应该在向量空间中出现在与地名词（如 France、Japan）相近的位置。
\end{knowledgebox}

\begin{figure}[H]
\centering
\includegraphics[width=0.95\textwidth]{slides-images/slide-040.jpg}
\caption{Recursive vs. Recurrent：递归网络沿树结构计算短语表示，循环网络沿序列计算\protect\footnotemark}
\end{figure}
\footnotetext{Slides 第41页。}

\subsection{简单递归神经网络}

最基本的树递归神经网络使用一个共享的权重矩阵 $W$ 在树的所有节点处进行组合：

$$\mathbf{p} = f\left(W \begin{bmatrix} \mathbf{c}_1 \\ \mathbf{c}_2 \end{bmatrix} + \mathbf{b}\right)$$

其中：
\begin{itemize}
    \item $\mathbf{c}_1, \mathbf{c}_2 \in \mathbb{R}^d$ 是两个子节点的向量表示
    \item $W \in \mathbb{R}^{d \times 2d}$ 是共享的组合矩阵
    \item $\mathbf{b}$ 是偏置向量
    \item $f$ 是非线性激活函数
    \item $\mathbf{p} \in \mathbb{R}^d$ 是父节点的表示
\end{itemize}

同时，可以学习一个评分函数来判断两个节点是否应该合并为一个成分：

$$\text{score} = \mathbf{u}^T \mathbf{p}$$

\begin{figure}[H]
\centering
\includegraphics[width=0.95\textwidth]{slides-images/slide-045.jpg}
\caption{贪心句法分析：从词向量开始，每一步选择得分最高的相邻对进行合并\protect\footnotemark}
\end{figure}
\footnotetext{Slides 第46页。}

\subsection{贪心分析算法}

利用上述评分函数，可以构建一个贪心的句法分析器：

\begin{enumerate}
    \item 初始化：所有词的词向量
    \item 计算所有相邻词对的合并得分
    \item 选择得分最高的对，合并为一个新节点
    \item 更新相邻关系，重新计算涉及新节点的得分
    \item 重复直到只剩一个根节点
\end{enumerate}

\begin{figure}[H]
\centering
\includegraphics[width=0.95\textwidth]{slides-images/slide-046.jpg}
\caption{贪心分析过程示例：``The cat sat on the mat''，首先合并``the cat''，然后``the mat''，逐步构建二叉分析树\protect\footnotemark}
\end{figure}
\footnotetext{Slides 第47页。}

\begin{warningbox}{简单 TreeRNN 的局限性}
使用单一权重矩阵 $W$ 在所有节点处组合存在明显局限：
\begin{itemize}
    \item 无法区分不同类型的语法组合（形容词+名词 vs 动词+宾语）
    \item 对所有词对采用完全相同的交互方式
    \item 表达能力受限，性能仅与 bigram Naive Bayes 分类器相当
\end{itemize}
\end{warningbox}

\subsection{本章小结}

树递归神经网络基于语言的递归结构和组合性原则，沿句法树自底向上构建短语和句子的向量表示。简单的 TreeRNN 使用共享权重矩阵进行节点组合，并可用于贪心句法分析。但单一权重矩阵的表达能力不足，需要更强大的组合函数。

%% ========== RNTN 与情感分析 ==========

